\documentclass[11pt,a4paper]{article}
\usepackage[T1]{fontenc}
\usepackage[utf8]{inputenc}
\usepackage{amsmath,amssymb,xcolor}
\usepackage[margin=25mm]{geometry}
\usepackage[colorlinks=true,urlcolor=blue]{hyperref}
\definecolor{nexorablue}{HTML}{2459B5}
\setlength{\parindent}{0pt}
\setlength{\parskip}{7pt}
\setlength{\emergencystretch}{2em}
\renewcommand{\familydefault}{\sfdefault}
\begin{document}
\makeatletter
\def\@oddfoot{\footnotesize Download for free at \url{https://openstax.org/details/books/university-physics-volume-2}.\hfil\thepage}
\let\@evenfoot\@oddfoot
\makeatother
{\small\bfseries\color{nexorablue} NEXORA CPGE}\par
{\small MPSI / Cours}\par
{\LARGE\bfseries Régimes transitoires}\par
\vspace{0.25em}\hrule\vspace{0.6em}

\section*{Objectifs et prérequis}
Revoir les lois de Kirchhoff, les conventions récepteur et les équations différentielles linéaires à coefficients constants. On suppose les dipôles idéaux et l’approximation des régimes quasi permanents applicable. Après une commutation, redessiner le circuit valable pour $t>0$.

\section*{Grandeurs d’état et continuité}
Dans la convention récepteur, $i_C=C\,\mathrm du_C/\mathrm dt$ et $u_L=L\,\mathrm di_L/\mathrm dt$. Intégrer autour de la commutation. Si le courant de la capacité et la tension de la bobine restent localement bornés, alors $u_C(0^+)=u_C(0^-)$ et $i_L(0^+)=i_L(0^-)$.

Pour déterminer les valeurs à gauche en régime continu établi, remplacer le condensateur par un circuit ouvert et la bobine idéale par un court-circuit. Réutiliser ces modèles pour obtenir les limites à droite si le nouveau régime permanent existe.

\section*{Circuit RC série}
Dans une boucle source constante $E$, résistance $R$ et condensateur $C$, écrire $E=Ri+u_C$, puis substituer $i=C\dot u_C$. L’écart à la valeur finale vérifie une équation homogène.
\[ RC\dot u_C+u_C=E,\qquad u_C(t)=E+(u_C(0^+)-E)e^{-t/(RC)}. \]
Pour une décharge dans $R$, prendre $E=0$. Pour une charge depuis zéro, $i(t)=(E/R)e^{-t/(RC)}$. À une constante de temps $RC$, environ $63{,}2\,\%$ du saut de tension est accompli.

\section*{Circuit RL série}
Écrire la loi des mailles en substituant la relation de la bobine. L’intensité tend vers la valeur donnée par Ohm et la tension de la bobine idéale tend vers zéro.
\[ L\dot i+Ri=E,\qquad i(t)=\frac ER+\left(i(0^+)-\frac ER\right)e^{-tR/L}. \]
La constante de temps est $L/R$. Lorsqu’une source est retirée mais que la boucle reste fermée par la résistance, l’intensité décroît sans sauter à zéro.

\section*{Circuit RLC série et trois régimes}
Avec la charge $q$, on obtient $L\ddot q+R\dot q+q/C=E$ et $\dot q=i$. Poser $x=q-CE$, $\alpha=R/(2L)$ et $\omega_0=1/\sqrt{LC}$. L’équation pour $x$ a pour polynôme caractéristique $r^2+2\alpha r+\omega_0^2$.

Lorsque $\alpha<\omega_0$, la réponse est $e^{-\alpha t}(A\cos\omega t+B\sin\omega t)$ avec $\omega=\sqrt{\omega_0^2-\alpha^2}$. Si $\alpha=\omega_0$, elle est $(A+Bt)e^{-\alpha t}$. Si $\alpha>\omega_0$, elle est une somme de deux exponentielles réelles décroissantes. Déterminer les deux constantes avec $q(0^+)$ et $i(0^+)$.

\section*{Bilan d’énergie}
Multiplier l’équation de maille par le courant. L’énergie emmagasinée est la somme des termes électrique et magnétique, tandis que la résistance dissipe de la puissance.
\[ \mathcal U=\frac{q^2}{2C}+\frac12Li^2,\qquad \frac{\mathrm d\mathcal U}{\mathrm dt}=Ei-Ri^2. \]
Pour une décharge sans source, l’énergie décroît même si la tension ou le courant oscille en changeant de signe.

\section*{Exemple et méthode CNC}
Avec $R=2\,\mathrm{k\Omega}$, $C=100\,\mathrm{\mu F}$ et une source de $5\,\mathrm V$, un condensateur initialement vide atteint $4\,\mathrm V$ à $t=RC\ln 5\simeq0{,}322\,\mathrm s$. La formule de seuil n’a de solution finie que si le seuil est strictement entre la tension initiale et la limite.

Avant tout calcul, distinguer la grandeur qui reste continue de la grandeur dont on cherche l’évolution. Vérifier l’unité de la constante de temps et la pente initiale à l’aide de la loi de maille.

\section*{Réduire le circuit vu par la capacité}
Dans un réseau linéaire de résistances et de sources continues, remplacer la partie connectée à la capacité par le générateur de Thévenin de tension $E_{\mathrm{th}}$ et résistance $R_{\mathrm{th}}$. Calculer la tension à vide, puis la résistance vue aux bornes avec les sources indépendantes annulées.

Pendant une phase où le circuit est fixe, la tension s’écrit $E_{\mathrm{th}}+(u_0-E_{\mathrm{th}})e^{-t/(R_{\mathrm{th}}C)}$. Une nouvelle commutation peut modifier les deux paramètres. Refaire le schéma équivalent pour chaque phase.

\section*{Commutations successives et changement d’origine}
Quand la seconde phase débute à $t=t_1$, utiliser le temps local $t-t_1$ et transmettre à cette phase la valeur du condensateur ou de la bobine juste avant la commutation. Pour une variable du premier ordre $x$, une phase écrit
\[ x(t)=x_\infty+(x(t_1^+)-x_\infty)e^{-(t-t_1)/\tau}. \]
Dans un RLC série, transmettre à la fois la charge et le courant si les tensions et courants concernés restent bornés. La pente peut changer brusquement lorsque l’excitation ou la résistance change.

\section*{Amortissement critique et identification sur une courbe}
Le seuil critique du RLC série est $R_c=2\sqrt{L/C}$. Pour une décroissance du premier ordre, la tangente initiale atteint la valeur finale au bout d’une constante de temps, et l’écart au régime permanent est réduit d’un facteur $e$ à cet instant.

Si plusieurs extrema oscillants sont visibles, estimer la pseudo-période entre deux maxima du même signe et l’amortissement à partir du rapport de leurs écarts à la valeur finale. Vérifier que la droite expérimentale de $\ln|x-x_\infty|$ est cohérente avec un modèle exponentiel seulement quand un tel modèle est justifié.

\section*{Énergie fournie et effet de la résistance}
Pour une charge RC depuis zéro sous une source idéale $E$, l’énergie fournie vaut $\int_0^\infty Ei\,\mathrm dt=E\int_0^{CE}\mathrm dq=CE^2$. La capacité en stocke $CE^2/2$ et la résistance dissipe l’autre moitié si $R>0$. Modifier $R$ change la vitesse de charge et la puissance maximale, mais pas ce résultat dans le modèle idéal.

Pour vérifier un calcul de puissance, distinguer le signe de la puissance reçue par chaque dipôle selon les conventions choisies.

\section*{Fiche de révision}
Trouver les états initial et final, poser l’écart à la limite, résoudre le système homogène et vérifier les signes. Au premier ordre, l’écart est multiplié par $e^{-1}$ après une constante de temps. Au second ordre, examiner le discriminant avant d’identifier le régime.

\bigskip
{\footnotesize Adaptation française par Nexora CPGE. Source : Samuel J. Ling, Jeff Sanny et William Moebs (OpenStax, Rice University),
\href{https://d3bxy9euw4e147.cloudfront.net/oscms-prod/media/documents/UniversityPhysicsVolume2-OP.pdf\#page=479}{\textit{University Physics Volume 2, édition 2018}}.
\href{https://creativecommons.org/licenses/by/4.0/}{CC BY 4.0}.}
\end{document}
